% !TEX root = ../main.tex

\chapter{数值方法}\label{ch:methodology}
  考虑方程 \eqref{eq:ppde} 的终值问题, 即
  \begin{equation}
    \begin{gathered}
      \begin{aligned}
        \frac{\partial u}{\partial t}(t, x)
        &+ \frac{1}{2} \operatorname{Tr}\!\big(
        \sigma\sigma^{\operatorname{T}}(t, x)\, \nabla_{x}^{2} u(t, x) \big)
         + \langle \mu(t, x), \nabla u(t, x) \rangle \\
        &+ f\!\big( t, x, u(t,x), \sigma^{\operatorname{T}}(t, x) \nabla u(t, x) \big)
        = 0,
        \quad (t,x) \in [0,T) \times \mathbb{R}^{d},
      \end{aligned} \\
      u(T, x) = g(x), \quad x \in \mathbb{R}^{d}.
    \end{gathered}
    \label{eq:ppdeg}
  \end{equation}
  为了应用 Itô 公式, 假设 $u \in C^{1, 2}([0, T) \times \mathbb{R}^{d})$,
  并要求 $\mu, \sigma, f, g$ 满足足够的正则性与增长条件.
  令状态过程满足
  \begin{equation}
    X_{t} = \xi + \int_{0}^{t} \mu(s, X_{s}) \dif s + \int_{0}^{t} \sigma(s, X_{s}) \dif W_{s},
    \label{eq:xt}
  \end{equation}
  再定义
  \begin{equation}
    Y_{t} \coloneqq u(t, X_{t}),
    \qquad
    Z_{t} \coloneqq \sigma^{\operatorname{T}}(t, X_{t}) \nabla u(t, X_{t}).
    \label{eq:ytzt}
  \end{equation}
  则原方程可写成 BSDE
  \begin{equation}
    \dif Y_{t} = -f(t, X_{t}, Y_{t}, Z_{t}) \dif t
    + Z_{t}^{\operatorname{T}} \dif W_{t},
    \label{eq:bsde}
  \end{equation}
  并满足终端条件 $Y_{T}=g(X_{T})$.
  从 \eqref{eq:ppdeg} 到 \eqref{eq:bsde} 的推导见附录~\ref{ch:derivation} 第\ref{sec:pde2bsde}节.

  取时间剖分 $0 = t_{0} < t_{1} < \cdots < t_{N} = T$, 其中 $\Delta t = T/N$,
  并记 $\Delta W_{k} = W_{t_{k + 1}} - W_{t_{k}}$.
  对 \eqref{eq:bsde} 采用 Euler 离散后, 得到
  \begin{equation}
    Y_{k + 1} = Y_{k} - \Delta t\, f\!(t_{k}, X_{k}, Y_{k}, Z_{k}) + \langle Z_{k}, \Delta W_{k} \rangle.
    \label{eq:dbsde}
  \end{equation}
  因而问题转化为: 在离散路径上同时近似初值 $Y_{0}$ 和各时间层的梯度项 $Z_{k}$,
  使终端值与 $g(X_{T})$ 尽可能一致.

  \section{线性方程的求解}\label{sec:linear}
    先考虑线性生成元
    \begin{equation}
      f\!( t, \boldsymbol{x}, y, \boldsymbol{z})
      = \mathcal{L}(t, \boldsymbol{x}) y
      + \langle \mathcal{M}(t, \boldsymbol{x}), \boldsymbol{z} \rangle
      + \mathcal{N}(t, \boldsymbol{x}).
      \label{eq:linearf}
    \end{equation}
    其中 $\langle \mathcal{M}, \boldsymbol{z} \rangle$ 可以并入漂移项处理,
    因此后文只考虑 $\mathcal{M}\equiv 0$ 的情形.

    对每条样本路径 $i$, 用随机特征映射逼近梯度项
    \begin{equation}
      Z_{k}^{(i)} \approx \alpha H_{k}^{(i)},
      \qquad
      H_{k}^{(i)} = \tanh(A x_{k}^{(i)} + b t_{k} + c),
      \label{eq:rff@l}
    \end{equation}
    其中 $A, b, c$ 为固定随机参数, $\alpha$ 为待训练的输出层参数.
    代入 \eqref{eq:dbsde} 后得到离散迭代
    \begin{equation}
      y_{k + 1}^{(i)}
      = \big( 1 - \Delta t\, \mathcal{L}(t_{k}, x_{k}^{(i)}) \big) y_{k}^{(i)}
      + (\Delta W_{k}^{(i)})^{\operatorname{T}} \alpha H_{k}^{(i)}
      - \Delta t\, \mathcal{N}(t_{k}, x_{k}^{(i)}).
      \label{eq:ldbsde}
    \end{equation}
    由于上式关于 $y_{0}$ 和 $\alpha$ 是线性的, 终端值 $y_{N}^{(i)}$ 也必然是它们的仿射函数.
    将递推完全展开后, 可写为
    \begin{equation}
      y_{N}^{(i)} = \beta_{i} y_{0} + \langle \Gamma_{i}, \alpha \rangle + \eta_{i},
      \label{eq:yN4i}
    \end{equation}
    其中 $\beta_{i}$, $\Gamma_{i}$, $\eta_{i}$ 由路径 $\{x_{k}^{(i)}, \Delta W_{k}^{(i)}\}$ 唯一确定,
    具体表达式分别见附录~\ref{ch:derivation} 第\ref{sec:linearder}节的
    \eqref{eq:beta}, \eqref{eq:gamma} 和 \eqref{eq:eta}.

    将所有路径的终端值拼接后, 便得到线性最小二乘问题.
    记
    \begin{equation}
      \vartheta =
      \begin{bmatrix}
        y_0 \\
        \operatorname{vec}(\alpha)
      \end{bmatrix},
      \qquad
      A_{i,\cdot} =
      \begin{bmatrix}
        \beta_i & \operatorname{vec}(\Gamma_i)^{\operatorname{T}}
      \end{bmatrix},
      \qquad
      B_i = g(x_N^{(i)})-\eta_i.
      \label{eq:linear_design}
    \end{equation}
    则线性方程最终求解
    \begin{equation}
      \min_{\vartheta}
      \lVert A\vartheta-B\rVert_2^2+\lambda_{\mathrm{r}}\lVert R\vartheta\rVert_2^2,
      \label{eq:llstsq}
    \end{equation}
    其中 $\lambda_{\mathrm{r}}>0$ 为 ridge 正则化参数,
    \begin{equation}
      R =
      \begin{bmatrix}
        0 & 0 \\
        0 & I_{dH}
      \end{bmatrix}.
      \label{eq:linpen}
    \end{equation}
    也就是说, 正则化只作用在随机特征输出层参数 $\alpha$ 上, 不惩罚初值 $y_{0}$.
    这样做可以避免把 $y_{0}$ 人为向零偏移, 同时抑制病态线性系统给出的过大 $\alpha$.

    实际计算时, 设未知量个数为 $p=1 + dH$.
    当样本数 $S \ge p$ 时, 直接求解 primal ridge 系统
    \begin{equation}
      \left(A^{\operatorname{T}}A+\lambda_{\mathrm{r}}R^{\operatorname{T}}R\right)\vartheta
      = A^{\operatorname{T}}B.
      \label{eq:primal}
    \end{equation}
    当 $S < p$ 时, 使用对偶形式降低线性系统规模.
    由于当前实现的对偶形式主要用于欠定情形, 其正则化写为
    \begin{equation}
      \vartheta =
      A^{\operatorname{T}} \left(AA^{\operatorname{T}}+\lambda_{\mathrm{r}}I_S\right)^{-1}B.
      \label{eq:dual}
    \end{equation}
    线性情形的优势在于不需要迭代更新随机特征参数,
    每次实验只需求解一个最小二乘系统即可.

  \section{半线性方程的求解}\label{sec:nonlinear}
    对一般半线性生成元 $f(t, \boldsymbol{x}, y, \boldsymbol{z})$,
    仍采用相同的随机特征近似
    \begin{equation}
      Z_{k}^{(i)} \approx \alpha H_k^{(i)}.
      \label{eq:rff@sl}
    \end{equation}
    此时离散迭代式为
    \begin{equation}
      y_{k + 1}^{(i)} = y_{k}^{(i)}
      - f\!(t_{k}, x_{k}^{(i)}, y_{k}^{(i)}, \alpha H_{k}^{(i)}) \Delta t
      + (\Delta W_{k}^{(i)})^{\operatorname{T}} \alpha H_{k}^{(i)}.
      \label{eq:nldbsde}
    \end{equation}
    由于生成元 $f$ 关于 $y$ 或 $z$ 不再保持线性, 终端值 $y_N^{(i)}$ 不能再写成关于参数的线性函数,
    因此末端匹配需要改写为半线性最小二乘问题.

    令参数向量为
    \begin{equation}
      \theta =
      \begin{bmatrix}
        y_{0} \\
        \operatorname{vec}(\alpha)
      \end{bmatrix},
      \label{eq:theta}
    \end{equation}
    对第 $i$ 条路径定义终端残差
    \begin{equation}
      r_{i}(\theta) = y_{N}^{(i)}(\theta) - g(x_{N}^{(i)}),
      \label{eq:residual4i}
    \end{equation}
    并记
    \begin{equation}
      r(\theta) =
      \begin{bmatrix}
        r_{1}(\theta) \\
        \vdots \\
        r_{S}(\theta)
      \end{bmatrix}.
      \label{eq:residual}
    \end{equation}
    最终求解问题为
    \begin{equation}
      \min_{\theta}
      \frac{1}{2} \lVert r(\theta)\rVert_{2}^{2}.
      \label{eq:nlls}
    \end{equation}

    本文采用 Levenberg--Marquardt 方法求解 \eqref{eq:nlls}.
    在第 $m$ 次迭代中, 将残差在当前参数 $\theta^{(m)}$ 处线性化:
    \begin{equation}
      r(\theta^{(m)}+\delta) \approx r(\theta^{(m)}) + J(\theta^{(m)})\delta,
      \label{eq:reslin}
    \end{equation}
    其中
    \begin{equation}
      J(\theta^{(m)}) = \frac{\partial r}{\partial \theta}(\theta^{(m)})
      \label{eq:jacobian}
    \end{equation}
    为雅可比矩阵.
    于是每步需要求解阻尼线性系统
    \begin{equation}
      (J^{\operatorname{T}}J + \lambda I)\delta^{(m)} = - J^{\operatorname{T}}r, \qquad \lambda > 0,
      \label{eq:lm}
    \end{equation}
    并更新
    \begin{equation}
      \theta^{(m + 1)} = \theta^{(m)} + \delta^{(m)}.
      \label{eq:theta+}
    \end{equation}
    若目标函数下降, 则接受该步并减小阻尼参数; 否则拒绝该步并增大阻尼参数.

    为了构造 \eqref{eq:lm}, 需要计算终端残差对参数的灵敏度.
    对第 $i$ 条路径定义
    \begin{equation}
      S_{k}^{(i)} = \frac{\partial y_{k}^{(i)}}{\partial \theta},
      \label{eq:Sk}
    \end{equation}
    则有初值
    \begin{equation}
      S_{0}^{(i)} =
      \begin{bmatrix}
        1 & 0 & \cdots & 0
      \end{bmatrix},
      \label{eq:S0}
    \end{equation}
    并满足递推
    \begin{equation}
      S_{k+1}^{(i)}
      =
      \left(1-\Delta t\, f_{y,k}^{(i)}\right)S_k^{(i)}
      +
      \left((\Delta W_k^{(i)})^{\operatorname{T}}
        -\Delta t\, f_{z,k}^{(i)}\right)
      \begin{bmatrix}
        0_{d\times 1} & I_d \otimes (H_k^{(i)})^{\operatorname{T}}
      \end{bmatrix},
      \label{eq:srecurrence}
    \end{equation}
    其中
    \begin{equation}
      f_{y,k}^{(i)} = \frac{\partial f}{\partial y} (t_{k}, x_{k}^{(i)}, y_{k}^{(i)}, \alpha H_{k}^{(i)}),
      \qquad
      f_{z,k}^{(i)} = \frac{\partial f}{\partial z} (t_{k}, x_{k}^{(i)}, y_{k}^{(i)}, \alpha H_{k}^{(i)}).
      \label{eq:fyfz}
    \end{equation}
    因而第 $i$ 条路径对应的雅可比行向量就是
    \begin{equation}
      J_{i}(\theta) = S_{N}^{(i)}
      = (\prod_{j = 0}^{N - 1} A_{j}^{(i)}) S_{0}^{(i)} +
      \sum_{k=0}^{N-1} ( \prod_{j =k + 1}^{N - 1} A_j^{(i)}) B_{k}^{(i)}.
      \label{eq:jacobian4i}
    \end{equation}
    上述灵敏度递推及其由来见附录~\ref{ch:derivation} 第\ref{sec:semilinearder}节.

    这与线性方程中系数矩阵的推导形式类似,
    区别在于此处 $A_{k}^{(i)}$ 和 $B_{k}^{(i)}$ 依赖当前迭代步中的 $y_{k}^{(i)}$ 和 $\alpha$,
    因而需要在每一次迭代中重新计算.
    但随机特征结构使待优化参数始终只出现在输出层,
    从而显著降低了优化问题的规模.
